\problemname{Lotteri}

På en karneval i din stad finns ett lotteri.
För att spela tar du en handfull lappar från en skål, där det står ett tal på varje lapp.
Det finns en exakt $50\%$ sannolikhet att du drar varje lapp.
Denna sannolikheten är oberoende för alla lappar.
Du vinner om summan av lapparnas tal du drar är exakt $T$.
Vad är sannolikheten att du vinner?

\section*{Indata}
Den första raden av indata innehåller heltalen $N$ och $T$ ($1 \leq N \leq 20$, $1 \leq T \leq 10^9$),
antalet lappar och det vinnande talet.

Nästa rad innehåller $N$ heltal $w_1, w_2, \dots, w_N$ ($1 \leq w_i \leq 5 \cdot 10^7$),
som beskriver talen på lapparna.

\section*{Utdata}
Skriv ut ett flyttal: sannolikheten att du vinner spelet om du spelar en gång, angiven i procent.
Svaret accepteras om skillnaden mellan ditt svar och domarens svar inte är mer än $10^{-4}$.
